% =====================================================================
%  1.1 Potenser — fasit   (MØNSTERSEKSJON)
%  \fk{overskrift}   \fo{oppgavenummer} svar
% =====================================================================

\fk{Del 1 — Potensreglene}
\fo{1.1} a) $3^6$ \ b) $5^4$ \ c) $2^{12}$ \ d) $a^6$
\fo{1.2} a) $1000$ \ b) $\frac{4}{9}$ \ c) $25$ \ d) $-25$
\fo{1.3} a) $a^7$ \ b) $a^5$
\fo{1.4} a) $27x^6$ \ b) $3a^3b$ \ c) $4x^2$ \ d) $x^6y$
\fo{1.5} a) $\frac{2^6}{2^5} = 2$ \ b) $\frac{3^6}{3^5} = 3$ \ c) $\frac{2^{10}}{2^9} = 2$
\fo{1.6} Feilen er overgangen til $\frac{8b^{18}}{4b^6}$: $b^9\cdot b^2=b^{11}$, ikke $b^{18}$. Riktig svar er $\frac{8b^{11}}{4b^6}=2b^5$.
\fo{1.7} a) Noen ganger. Venstre side er $a^5$, og $a^5 = a^6$ bare når $a = 0$ eller $a = 1$.
b) Alltid. $(-a)^3 = (-1)^3\cdot a^3 = -a^3$.
c) Noen ganger. Venstre side er $a^4$ ($a\neq0$), og $a^4=a^3$ bare når $a=1$. Verdien $a=0$ er ikke tillatt.
d) Aldri. Venstre side er $2\cdot3^n$, høyre side er $3\cdot3^n$, og $2\cdot3^n\neq3\cdot3^n$. (Derimot stemmer $2^n+2^n=2^{n+1}$ alltid.)
\fo{1.8} $2^{12}=(2^2)^6=4^6$, $2^{12}=(2^3)^4=8^4$ og $2^{12}=(2^4)^3=16^3$. Med 32 må $32^x=(2^5)^x=2^{5x}=2^{12}$, så $5x=12$ og $x=\frac{12}{5}$. Det er ikke et helt tall.
\fo{Sjekk 1} a) $x^7$ \ b) $9$ og $-9$ \ c) $(2\cdot2\cdot2)\cdot(2\cdot2\cdot2\cdot2)$ er sju faktorer 2.

\fk{Del 2 — Null, negative eksponenter og standardform}
\fo{2.1} a) $1$ \ b) $\frac{1}{36}$ \ c) $\frac{1}{100} = 0,01$ \ d) $2$
\fo{2.2} a) $5,2 \cdot 10^4$ \ b) $7,1 \cdot 10^{-3}$ \ c) $3,4 \cdot 10^7$
\fo{2.3} a) $\frac{1}{x^2}$ \ b) $\frac{3}{a}$ \ c) $b^4$
\fo{2.4} a) $a$ \ b) $\frac{8}{x^3}$ \ c) $\frac{2x^3}{y^3}$ \ d) $a^2b^4$
\fo{2.5} a) $9 \cdot 10^{-4}$ \ b) $4 \cdot 10^6$ \ c) $27 \cdot 10^{-6} = 2,7 \cdot 10^{-5}$
\fo{2.6} $8,3 \cdot 60\text{ s} = 498\text{ s}$. Avstand: $498 \cdot 3,0 \cdot 10^8\text{ m} \approx 1,5 \cdot 10^{11}\text{ m} = 1,5 \cdot 10^8\text{ km}$.
\fo{2.7} a) Aldri. $2^{-n}=\frac{1}{2^n}$ er alltid positivt.
b) Noen ganger. Med $a=2$ er $\frac12<2$, sant. Med $a=\frac12$ er $2<\frac12$, usant. Sant for $a>1$ og $-1<a<0$ ($a\neq0$).
c) Noen ganger. $(2\cdot10^3)(3\cdot10^4)=6\cdot10^7$ er standardform, men $(5\cdot10^3)(4\cdot10^2)=20\cdot10^5$ er ikke det.
\fo{2.8} $2^{-10} = \frac{1}{1024}$, $10^{-3} = \frac{1}{1000}$ og $4^{-4} = \frac{1}{256}$.
Minst nevner gir størst brøk, så $4^{-4}$ er størst.
\fo{Sjekk 2} a) $\frac{1}{9}$ \ b) $4,5 \cdot 10^{-4}$ \ c) Nei. $25 > 10$. Standardform: $2,5 \cdot 10^4$.

\fk{Ekstra utfordring}
\fo{U1} Teller: $8a^6b^{-3} \cdot a^2b^2 = 8a^8b^{-1}$. Brøken blir $\frac{8a^8b^{-1}}{4a^7b^{-2}} = 2ab$.
\fo{U2} $\frac{2^5 \cdot 2^3 \cdot 3^3}{3^2 \cdot 2^4 \cdot 3^2} = \frac{2^8 \cdot 3^3}{2^4 \cdot 3^4} = \frac{2^4}{3} = \frac{16}{3}$
\fo{U3} a) $2^x = 2^{-5}$, $x = -5$ \ b) $3^{2x} = 3^3$, $x = \frac{3}{2}$ \ c) $2^{2x+2} = 2^{3x}$, $x = 2$
\fo{U4} a) $20$. Programmet øker $n$ så lenge $2^n < 10^6$. $2^{19} = 524\,288$ og $2^{20} = 1\,048\,576$.
b) Bytt linje 2 med \lstinline[style=laereverk,numbers=none,frame=none]|while 3**n <= 10**9:|. Svaret blir $19$.
\fo{U5} $2^{100} = (2^5)^{20} = 32^{20}$ og $3^{60} = (3^3)^{20} = 27^{20}$. Siden $32 > 27$, er $2^{100}$ størst.
\fo{U6} $4^6 \cdot 5^{10} = 2^{12} \cdot 5^{10} = 2^2 \cdot (2 \cdot 5)^{10} = 4 \cdot 10^{10}$. Tallet har 11 sifre.
\fo{U7} a) $m = 3$, $n = 2$: $(a^3)^2 = a^6$, men $a^{(3^2)} = a^9$.
b) Vi trenger $m \cdot n = m^n$. Når $n = 1$, stemmer det for alle $m$. Når $n = 2$, får vi $2m = m^2$, så $m = 2$.
Når $n \ge 3$, er $m^n \ge 2^{n-1} \cdot m > n \cdot m$ for $m \ge 2$, og $m = 1$ gir $1 \ne n$.
Svar: $n = 1$ (alle $m$), eller $m = n = 2$.
